% =========================================================
% Nowhere-Dense Sets and a Preview of Baire's Theorem
% Source: volume-iii/analysis/metric-spaces/notes/metric-spaces/
%
% Dependencies:
%   notes-metric-space.tex       — def:metric
%   notes-open-closed-sets.tex   — def:open, def:closed, def:interior
%   notes-dense-separable.tex    — def:dense
%   notes-completeness.tex       — completeness
% =========================================================

% ---------------------------------------------------------
\section{Nowhere-Dense Sets and Baire's Theorem}
\label{sec:nowhere-dense-baire}
% ---------------------------------------------------------

% ---- Toolkit ----
\begin{tcolorbox}[colback=gray!6, colframe=gray!40, arc=2pt,
  left=6pt, right=6pt, top=4pt, bottom=4pt,
  title={\small\textbf{Nowhere-Dense Sets and Baire --- Quick Reference}},
  fonttitle=\small\bfseries]
\small
\begin{tabular}{l l l l}
\toprule
\textbf{Label} & \textbf{Statement} & \textbf{Proof method} & \textbf{Detail} \\
\midrule
D\ref{def:nowhere-dense}
  & $E$ nowhere-dense iff $\operatorname{int}(\overline{E}) = \varnothing$
  & Definition
  & \hyperref[def:nowhere-dense]{$\downarrow$} \\[4pt]
D\ref{def:meager}
  & Meager $=$ countable union of nowhere-dense sets
  & Definition
  & \hyperref[def:meager]{$\downarrow$} \\[4pt]
P\ref{prop:nowhere-dense-equiv}
  & $E$ nowhere-dense $\iff$ $X \setminus \overline{E}$ is dense
  & Complement characterization
  & \hyperref[prop:nowhere-dense-equiv]{$\downarrow$} \\[4pt]
T\ref{thm:baire-preview}
  & Baire: complete metric space is not meager in itself
  & Statement only (proof deferred)
  & \hyperref[thm:baire-preview]{$\downarrow$} \\[4pt]
\bottomrule
\end{tabular}
\end{tcolorbox}

% ---------------------------------------------------------
\paragraph{Motivation.}
Density gives us a notion of ``large'' subsets: a dense set reaches
every corner of the space.
But we also need a notion of ``small'' or ``thin'' subsets,
one that captures sets which --- no matter how you zoom in ---
never fill any open region.

The integers $\mathbb{Z}$ in $\mathbb{R}$ are a perfect example.
No matter which open interval you look at, $\mathbb{Z}$ does not
fill it; $\mathbb{Z}$ has gaps everywhere.
This is the idea behind \emph{nowhere-dense} sets.

Understanding thin sets leads to one of the deepest structural results
in metric space theory: Baire's theorem, which says that a complete
metric space cannot be built from countably many thin pieces.
We introduce the definitions here and state the theorem;
the proof belongs to the compactness chapter where the key tools are available.

% ---- Definition: Nowhere-Dense ----
\begin{tcolorbox}[colback=defbox, colframe=defborder, arc=2pt,
  left=6pt, right=6pt, top=4pt, bottom=4pt,
  title={\small\textbf{Definition (Nowhere-Dense)}},
  fonttitle=\small\bfseries]
\begin{definition}[Nowhere-Dense]
\label{def:nowhere-dense}
Let $(X,d)$ be a metric space and $E \subseteq X$.
We say $E$ is \emph{nowhere-dense in $X$} if the interior of its closure
is empty:
\[
\operatorname{int}(\overline{E}) = \varnothing.
\]
\end{definition}
\end{tcolorbox}

\begin{remark*}[Fully quantified form]
\[
E \text{ is nowhere-dense}
\;\iff\;
\forall\, x \in \overline{E},\; \forall\, r > 0,\;
B_r(x) \not\subseteq \overline{E}.
\]
No point of $\overline{E}$ has an open ball entirely inside $\overline{E}$.
\end{remark*}

\begin{remark*}[Intuition: nowhere-dense vs.\ not-dense]
These sound similar but are very different.

A set $E$ is \emph{not dense} if there is some open ball it misses entirely:
$\exists\, B_r(x)$ with $B_r(x) \cap E = \varnothing$.

A set $E$ is \emph{nowhere-dense} if \emph{every} open ball contains a
sub-ball that $E$ misses entirely.
Nowhere-dense is a much stronger condition --- it says $E$ fails to be
dense even when restricted to any open region.

Think of nowhere-dense as: ``thin no matter where you zoom in.''
Not-dense just means ``thin somewhere.''
\end{remark*}

\begin{remark*}[Why close before taking interior?]
We use $\operatorname{int}(\overline{E})$ rather than $\operatorname{int}(E)$
because isolated points would otherwise cause confusion.
A single point $\{p\}$ has empty interior but its closure is still $\{p\}$,
which also has empty interior --- so $\{p\}$ is nowhere-dense, as it should be.
The closure step ensures the definition behaves well for any set,
not just closed ones.
For closed sets $E$, the definition simplifies to
$\operatorname{int}(E) = \varnothing$.
\end{remark*}

% ---- Proposition: equivalent characterization ----

\begin{proposition}[Complement characterization of nowhere-dense]
\label{prop:nowhere-dense-equiv}
$E \subseteq X$ is nowhere-dense if and only if $X \setminus \overline{E}$
is dense in $X$.
\end{proposition}

\begin{proof}
By definition, $E$ is nowhere-dense iff $\operatorname{int}(\overline{E}) = \varnothing$.
The interior of $\overline{E}$ is the largest open set contained in
$\overline{E}$, so $\operatorname{int}(\overline{E}) = \varnothing$ iff no
non-empty open set is contained in $\overline{E}$, iff every non-empty open
set meets $X \setminus \overline{E}$.
By Proposition~\ref{prop:dense-equiv}(iii), this is exactly
$\overline{X \setminus \overline{E}} = X$, i.e.\ $X \setminus \overline{E}$
is dense.
\end{proof}

\begin{remark*}[Reading the characterization]
$X \setminus \overline{E}$ dense says: the complement of $\overline{E}$
meets every open set.
In other words, no matter where you look, you can find points
that are not even close to $E$.
This is the precise form of ``$E$ is thin everywhere.''
\end{remark*}

% ---- Examples ----

\begin{example}[Nowhere-dense sets]
\label{ex:nowhere-dense}
\begin{enumerate}[label=(\alph*)]
  \item \textbf{Integers in $\mathbb{R}$.}
        $\mathbb{Z}$ is closed, so $\overline{\mathbb{Z}} = \mathbb{Z}$.
        The interior of $\mathbb{Z}$ is empty (no open interval is contained
        in $\mathbb{Z}$). So $\mathbb{Z}$ is nowhere-dense.

  \item \textbf{The Cantor set $C \subset [0,1]$.}
        $C$ is closed and has empty interior (it contains no intervals).
        So $C$ is nowhere-dense in $\mathbb{R}$,
        despite being uncountable --- a striking demonstration that
        nowhere-dense does not mean small in the sense of cardinality.

  \item \textbf{A single point $\{p\}$.}
        $\overline{\{p\}} = \{p\}$ and $\operatorname{int}(\{p\}) = \varnothing$.
        So every singleton is nowhere-dense.

  \item \textbf{A closed ball's boundary $\partial B_r(x)$ in $\mathbb{R}^n$.}
        The boundary sphere is closed with empty interior, hence nowhere-dense.
\end{enumerate}
\end{example}

\begin{example}[Sets that are NOT nowhere-dense]
\label{ex:not-nowhere-dense}
\begin{enumerate}[label=(\alph*)]
  \item \textbf{$\mathbb{Q}$ in $\mathbb{R}$.}
        $\overline{\mathbb{Q}} = \mathbb{R}$, and
        $\operatorname{int}(\mathbb{R}) = \mathbb{R} \neq \varnothing$.
        So $\mathbb{Q}$ is \emph{not} nowhere-dense,
        even though $\mathbb{Q}$ itself has empty interior.
        (The closure step is essential.)

  \item \textbf{Any open set $U \neq \varnothing$.}
        $\operatorname{int}(\overline{U}) \supseteq \operatorname{int}(U) = U \neq \varnothing$.
        No non-empty open set is nowhere-dense.
\end{enumerate}
\end{example}

% ---- Definition: Meager ----
\begin{tcolorbox}[colback=defbox, colframe=defborder, arc=2pt,
  left=6pt, right=6pt, top=4pt, bottom=4pt,
  title={\small\textbf{Definition (Meager / First Category)}},
  fonttitle=\small\bfseries]
\begin{definition}[Meager Set]
\label{def:meager}
A subset $A \subseteq X$ is \emph{meager} (or \emph{of first category})
if it can be written as a countable union of nowhere-dense sets:
\[
A = \bigcup_{n=1}^{\infty} E_n,
\quad \text{where each } E_n \text{ is nowhere-dense.}
\]
A set that is not meager is called \emph{non-meager} (or
\emph{of second category}).
\end{definition}
\end{tcolorbox}

\begin{remark*}[Intuition: meager as ``topologically small'']
Meager sets are thin in every topological sense.
A countable union of sets each of which is ``thin everywhere'' is
still considered ``thin.''
This gives us a topological analogue of measure-zero sets:
\begin{center}
\begin{tabular}{l l}
\toprule
\textbf{Measure theory} & \textbf{Baire category} \\
\midrule
Measure-zero set & Nowhere-dense set \\
Null set (countable union of measure-zero) & Meager set \\
Sets of positive measure & Non-meager (second category) sets \\
\bottomrule
\end{tabular}
\end{center}
The analogy is imperfect --- a set can be meager and have full measure,
or be non-meager and have measure zero --- but the intuition of
``negligible'' is the right one.
\end{remark*}

\begin{example}[$\mathbb{Q}$ is meager in $\mathbb{R}$]
\label{ex:Q-meager}
Write $\mathbb{Q} = \{q_1, q_2, q_3, \ldots\}$ (countable).
Each singleton $\{q_n\}$ is nowhere-dense.
So $\mathbb{Q} = \bigcup_{n=1}^\infty \{q_n\}$ is meager in $\mathbb{R}$.

This says that the rationals, despite being dense, are
\emph{topologically small}.
Density and meagerness are not opposites: $\mathbb{Q}$ is simultaneously
dense and meager in $\mathbb{R}$.
\end{example}

% ---- Baire's Theorem: Statement only ----

\begin{tcolorbox}[colback=thmbox, colframe=thmborder, arc=2pt,
  left=6pt, right=6pt, top=4pt, bottom=4pt,
  title={\small\textbf{Theorem (Baire Category Theorem --- Abbott \S8.2,
         Pugh Theorem 32)}},
  fonttitle=\small\bfseries]
\begin{theorem}[Baire Category Theorem]
\label{thm:baire-preview}
Let $(X, d)$ be a complete metric space.
Then $X$ is not meager in itself: $X$ cannot be written as a countable
union of nowhere-dense sets.

Equivalently, if $X = \bigcup_{n=1}^\infty E_n$ then at least one $E_n$
has non-empty interior.
\end{theorem}
\end{tcolorbox}

\begin{remark*}[Proof deferred]
The standard proof uses a diagonal/nested-ball argument that requires
completeness in an essential way.
The structure is: choose points from a sequence of shrinking open balls,
each forcing us deeper into an open dense set; completeness ensures the
resulting Cauchy sequence converges to a point inside all of them.
The proof will appear once the compactness tools are in place.
\end{remark*}

\begin{remark*}[What Baire's theorem says about $\mathbb{R}$]
Since $\mathbb{R}$ is complete, it is non-meager: you cannot cover $\mathbb{R}$
by countably many nowhere-dense sets.
But we just showed $\mathbb{Q}$ is meager.
Baire's theorem says the irrationals $\mathbb{R} \setminus \mathbb{Q}$
cannot be meager: if they were, then
$\mathbb{R} = \mathbb{Q} \cup (\mathbb{R} \setminus \mathbb{Q})$
would be the union of two meager sets, hence meager, contradicting
the theorem.

This is a precise sense in which the irrationals are
``much larger'' than the rationals, independent of cardinality:
the irrationals are non-meager, the rationals are meager.
\end{remark*}

\begin{remark*}[Application: most continuous functions are
  nowhere differentiable]
One of the most striking consequences of Baire's theorem is the
following result in $C([0,1])$ with the uniform metric.
Let $D \subseteq C([0,1])$ be the set of functions that are
differentiable at at least one point.
One can show that $D$ is meager in $C([0,1])$.
Since $C([0,1])$ is complete (and hence non-meager by Baire),
$D$ does not cover $C([0,1])$.
Therefore most continuous functions --- in the precise sense of
non-meager complement --- are differentiable \emph{nowhere}.
The differentiable functions, which seem like the ``usual'' ones
from calculus, are in fact the exceptional case.
\end{remark*}

\begin{remark*}[Completeness is essential]
The Baire Category Theorem is false for incomplete spaces.
The rationals $\mathbb{Q}$ (with the usual metric, which is incomplete)
\emph{can} be written as a countable union of singletons, each
nowhere-dense in $\mathbb{Q}$.
So $\mathbb{Q}$ is meager in itself.
Completeness is not just a convenient hypothesis --- it is what prevents
the space from being ``built from nothing.''
\end{remark*}
